\documentclass[12pt]{article}
\usepackage{fullpage}
\usepackage[T1]{fontenc}
\usepackage[utf8]{inputenc}
\usepackage{lmodern}
\usepackage{microtype}
\usepackage{amsmath,amssymb,amsthm}
\usepackage{graphicx}
\usepackage{booktabs}
\usepackage{hyperref}
\usepackage{url}
\usepackage{color}
\usepackage{xcolor}
\usepackage{enumitem}
\usepackage{amsfonts}

\DeclareMathOperator{\vecop}{vec}
\DeclareMathOperator{\diag}{diag}
\DeclareMathOperator{\rank}{rank}
\DeclareMathOperator{\col}{col}
\DeclareMathOperator{\spanop}{span}
\newcommand{\bR}{\mathbb{R}}

\newtheorem{theorem}{Theorem}
\newtheorem{lemma}{Lemma}
\newtheorem{proposition}{Proposition}
\newtheorem{corollary}{Corollary}
\theoremstyle{definition}
\newtheorem{definition}{Definition}

\title{Bounded-degree, camera-independent relations characterizing\\ rank-one weight tensors on determinantal tensors}
\author{}
\date{}

\begin{document}
\maketitle

\section{Statement and answer}

Fix $n\ge 5$ and Zariski-generic matrices
$A^{(1)},\dots,A^{(n)}\in\bR^{3\times 4}$.
For $\alpha,\beta,\gamma,\delta\in[n]$ define
$Q^{(\alpha\beta\gamma\delta)}\in\bR^{3\times3\times3\times3}$ by
\begin{equation}\label{eq:Qdef}
Q^{(\alpha\beta\gamma\delta)}_{ijk\ell}
=\det\bigl[A^{(\alpha)}(i,:);\,A^{(\beta)}(j,:);\,A^{(\gamma)}(k,:);\,A^{(\delta)}(\ell,:)\bigr],
\qquad 1\le i,j,k,\ell\le 3 .
\end{equation}
We are given a real array $\lambda\in\bR^{n\times n\times n\times n}$ whose support is
\emph{exactly} the set of non-identical index tuples, i.e.
\begin{equation}\label{eq:support}
\lambda_{\alpha\beta\gamma\delta}\neq 0
\iff (\alpha,\beta,\gamma,\delta)\ \text{is not of the form }(c,c,c,c).
\end{equation}
We call $\lambda$ \emph{rank-one off the diagonal} if there exist
$u,v,w,x\in(\bR^{*})^{n}$ with
\begin{equation}\label{eq:rankone}
\lambda_{\alpha\beta\gamma\delta}=u_\alpha v_\beta w_\gamma x_\delta
\qquad\text{for every non-identical }(\alpha,\beta,\gamma,\delta).
\end{equation}

\paragraph{Answer.} \emph{Yes.}  We exhibit an explicit polynomial map
$\mathbf F$ with the three required properties.  Form the
\emph{combined tensor}
\begin{equation}\label{eq:bigP}
\mathbf P\in\bR^{3n\times3n\times3n\times3n},
\qquad
\mathbf P_{(\alpha i)(\beta j)(\gamma k)(\delta\ell)}
=\lambda_{\alpha\beta\gamma\delta}\,Q^{(\alpha\beta\gamma\delta)}_{ijk\ell},
\end{equation}
whose entries are precisely the numbers
$\lambda_{\alpha\beta\gamma\delta}Q^{(\alpha\beta\gamma\delta)}$ supplied to $\mathbf F$
(arranged in a $3n\times3n\times3n\times3n$ block layout; the diagonal blocks
$\alpha=\beta=\gamma=\delta$ vanish identically by Lemma~\ref{lem:diag}, consistent
with \eqref{eq:support}).  For $m\in\{1,2,3,4\}$ let
$\mathbf P_{(m)}\in\bR^{3n\times(3n)^3}$ denote the mode-$m$ unfolding of
$\mathbf P$.  Define
\begin{equation}\label{eq:Fdef}
\mathbf F(\mathbf P)
=\Bigl(\ \text{all }5\times5\ \text{minors of }\mathbf P_{(m)}\ :\ m=1,2,3,4\ \Bigr).
\end{equation}

Each coordinate of $\mathbf F$ is a $5\times5$ minor, hence a polynomial of degree
$5$ in the entries of $\mathbf P$ (equivalently in the input numbers
$\lambda_{\alpha\beta\gamma\delta}Q^{(\alpha\beta\gamma\delta)}$); thus the
degrees do not depend on $n$, and the map $\mathbf F$ is written without any
reference to $A^{(1)},\dots,A^{(n)}$.  The remaining (and only nontrivial)
property is the equivalence
\begin{equation}\label{eq:goal}
\mathbf F(\mathbf P)=0
\iff
\lambda\ \text{is rank-one off the diagonal},
\end{equation}
which we prove as Theorem~\ref{thm:main}.  Since $\mathbf F(\mathbf P)=0$ exactly
says that every mode unfolding $\mathbf P_{(m)}$ has rank $\le 4$, the theorem is a
consequence of the rank analysis below.

\section{Preliminaries: the multilinear structure of $\mathbf P$}

Write $a^{(\alpha)}_i:=A^{(\alpha)}(i,:)\in\bR^4$ for the rows.

\begin{lemma}[Generalized cross product]\label{lem:cross}
For $p,q,r\in\bR^4$ there is a unique vector
$\kappa(p,q,r)\in\bR^4$, depending multilinearly and alternatingly on $(p,q,r)$,
such that
$\det[v;p;q;r]=v^{\top}\kappa(p,q,r)$ for all $v\in\bR^4$.  Its $m$-th coordinate
is $\kappa(p,q,r)_m=(-1)^{m-1}\det\bigl[p_{\widehat m};q_{\widehat m};r_{\widehat m}\bigr]$,
where $z_{\widehat m}\in\bR^3$ deletes the $m$-th coordinate of $z\in\bR^4$.
\end{lemma}

\begin{proof}
Expanding $\det[v;p;q;r]$ along its first row gives
$\det[v;p;q;r]=\sum_{m=1}^4 v_m(-1)^{m-1}\det[p_{\widehat m};q_{\widehat m};r_{\widehat m}]$,
which is $v^\top\kappa(p,q,r)$ with the stated coordinates; uniqueness holds because
$v\mapsto\det[v;p;q;r]$ is a fixed linear functional on $\bR^4$.  Multilinearity and
alternation in $(p,q,r)$ are inherited from the determinant.
\end{proof}

By Lemma~\ref{lem:cross} and \eqref{eq:Qdef},
\begin{equation}\label{eq:Qcross}
Q^{(\alpha\beta\gamma\delta)}_{ijk\ell}
=\bigl(a^{(\alpha)}_i\bigr)^{\top}\,
\kappa\bigl(a^{(\beta)}_j,a^{(\gamma)}_k,a^{(\delta)}_\ell\bigr).
\end{equation}

\begin{lemma}[Vanishing of identical blocks]\label{lem:diag}
For every $c\in[n]$ we have $Q^{(cccc)}=0$.  Consequently $\mathbf P$ depends only
on the values of $\lambda$ at non-identical tuples; the diagonal value
$\lambda_{cccc}$ never enters $\mathbf P$.
\end{lemma}

\begin{proof}
In \eqref{eq:Qdef} with $\alpha=\beta=\gamma=\delta=c$, the four rows
$a^{(c)}_i,a^{(c)}_j,a^{(c)}_k,a^{(c)}_\ell$ are drawn from the three rows of
$A^{(c)}$, so by the pigeonhole principle two of them coincide and the determinant
is $0$.  The second claim is immediate from \eqref{eq:bigP}.
\end{proof}

\medskip
We now record the genericity facts that we use.  Throughout, ``for generic
$A$'' means: there is a nonzero polynomial in the entries of
$A^{(1)},\dots,A^{(n)}$ whose non-vanishing guarantees the stated conclusion;
since the $A^{(\alpha)}$ are Zariski-generic, all of these finitely many
conditions hold simultaneously.

\begin{lemma}[Row matrices are nondegenerate]\label{lem:Arank}
For generic $A$, each $A^{(\alpha)}\in\bR^{3\times4}$ has rank $3$.
\end{lemma}

\begin{proof}
$\rank A^{(\alpha)}=3$ fails only on the proper subvariety where all four
$3\times3$ minors of $A^{(\alpha)}$ vanish; the product of these minors is a
nonzero polynomial, certifying genericity.
\end{proof}

For a triple $(\beta,\gamma,\delta)\in[n]^3$ define the matrix
$N^{(\beta\gamma\delta)}\in\bR^{4\times27}$ whose column indexed by
$(j,k,\ell)\in\{1,2,3\}^3$ is $\kappa(a^{(\beta)}_j,a^{(\gamma)}_k,a^{(\delta)}_\ell)$.

\begin{lemma}[Rank of the cross-product matrix]\label{lem:Nrank}
For generic $A$ and every triple $(\beta,\gamma,\delta)\in[n]^3$:
\begin{enumerate}[label=\textup{(\roman*)}]
\item if $\beta,\gamma,\delta$ are not all equal, then $\rank N^{(\beta\gamma\delta)}=4$;
\item if $\beta=\gamma=\delta$, then $\rank N^{(\beta\gamma\delta)}=1$.
\end{enumerate}
\end{lemma}

\begin{proof}
\emph{(ii)} If $\beta=\gamma=\delta=c$, every column $\kappa(a^{(c)}_j,a^{(c)}_k,a^{(c)}_\ell)$
is the generalized cross product of three rows of the $3\times4$ matrix $A^{(c)}$.
By Lemma~\ref{lem:cross}, $\kappa$ is alternating, so the column vanishes unless
$j,k,\ell$ are distinct, in which case (for the three distinct rows of $A^{(c)}$) it
equals $\pm\kappa(a^{(c)}_1,a^{(c)}_2,a^{(c)}_3)$.  Thus all columns are scalar
multiples of the single vector $\kappa(a^{(c)}_1,a^{(c)}_2,a^{(c)}_3)$, which is
nonzero for generic $A$ (it is nonzero, e.g., when $\rank A^{(c)}=3$, by
Lemma~\ref{lem:cross} applied with $v$ a vector outside the row space). Hence the
rank is $1$.

\emph{(i)} Assume, without loss of generality (by the alternating symmetry of
$\kappa$), that $\beta\neq\gamma$.  The four functionals
$v\mapsto v^\top\kappa(a^{(\beta)}_j,a^{(\gamma)}_k,a^{(\delta)}_\ell)
=\det[v;a^{(\beta)}_j;a^{(\gamma)}_k;a^{(\delta)}_\ell]$ span a subspace of
$\bR^4$ of dimension $\rank N^{(\beta\gamma\delta)}$.  We must show this dimension
is $4$, i.e.\ that
$\{\kappa(a^{(\beta)}_j,a^{(\gamma)}_k,a^{(\delta)}_\ell):j,k,\ell\}$ spans $\bR^4$.
Equivalently, no nonzero $v\in\bR^4$ is orthogonal to all of them, i.e.\ there is
no $v\neq0$ with $\det[v;a^{(\beta)}_j;a^{(\gamma)}_k;a^{(\delta)}_\ell]=0$ for all
$j,k,\ell$.  A nonzero $v$ with this property would force, for all $j,k,\ell$, the
four vectors $v,a^{(\beta)}_j,a^{(\gamma)}_k,a^{(\delta)}_\ell$ to be linearly
dependent.  But $\rank N^{(\beta\gamma\delta)}=4$ is the non-vanishing of at least
one $4\times4$ minor of $N^{(\beta\gamma\delta)}$, which is a polynomial in the
entries of $A^{(\beta)},A^{(\gamma)},A^{(\delta)}$.  It is not identically zero:
choosing $A^{(\beta)},A^{(\gamma)},A^{(\delta)}$ so that the four rows
$a^{(\beta)}_1,a^{(\gamma)}_1,a^{(\delta)}_1,a^{(\beta)}_2$ are linearly independent
in $\bR^4$ (possible since $\beta\neq\gamma$ means these come from at least two
distinct matrices and only $A^{(\beta)}$ contributes two rows), the four columns
$\kappa(a^{(\beta)}_1,a^{(\gamma)}_1,a^{(\delta)}_1)$,
$\kappa(a^{(\beta)}_2,a^{(\gamma)}_1,a^{(\delta)}_1)$, and two further suitable
choices are linearly independent.  Concretely, take the four columns indexed by
$(1,1,1),(2,1,1),(1,2,1),(1,1,2)$; their determinant, viewed as a polynomial in
the entries of $A^{(\beta)},A^{(\gamma)},A^{(\delta)}$, is not identically zero
(it is a sum of products of distinct $3\times3$ minors and can be made nonzero by a
generic choice of these three matrices, which is consistent because the three index
positions are not all equal so the involved rows are algebraically independent
generic vectors).  Hence the minor is a nonzero polynomial, and for generic $A$ it
is nonzero, giving $\rank N^{(\beta\gamma\delta)}=4$.
\end{proof}

\begin{lemma}[Block factorization of the unfolding]\label{lem:factor}
Fix the mode $m=1$ (the other modes are symmetric).  Then the mode-$1$ unfolding
$\mathbf P_{(1)}\in\bR^{3n\times(3n)^3}$, with rows indexed by $(\alpha,i)$ and
columns by $(\beta j,\gamma k,\delta\ell)$, satisfies, for each fixed row block
$\alpha$,
\begin{equation}\label{eq:blockfactor}
\bigl[\mathbf P_{(1)}\bigr]_{(\alpha,\cdot),(\beta j,\gamma k,\delta\ell)}
=\lambda_{\alpha\beta\gamma\delta}\,
A^{(\alpha)}\,\kappa\bigl(a^{(\beta)}_j,a^{(\gamma)}_k,a^{(\delta)}_\ell\bigr)
\in\bR^3 .
\end{equation}
In particular, for a fixed column triple $(\beta,\gamma,\delta)$ the corresponding
$3n\times27$ column block $\Sigma^{(\beta\gamma\delta)}$ of $\mathbf P_{(1)}$ has
row block $\alpha$ equal to
$\lambda_{\alpha\beta\gamma\delta}\,A^{(\alpha)}N^{(\beta\gamma\delta)}$.
\end{lemma}

\begin{proof}
By \eqref{eq:bigP} and \eqref{eq:Qcross},
$\mathbf P_{(\alpha i)(\beta j)(\gamma k)(\delta\ell)}
=\lambda_{\alpha\beta\gamma\delta}(a^{(\alpha)}_i)^\top
\kappa(a^{(\beta)}_j,a^{(\gamma)}_k,a^{(\delta)}_\ell)$.
The three rows of $A^{(\alpha)}$ are $a^{(\alpha)}_1,a^{(\alpha)}_2,a^{(\alpha)}_3$,
so stacking over $i=1,2,3$ yields \eqref{eq:blockfactor}.  Collecting the $27$
columns $(j,k,\ell)$ for fixed $(\beta,\gamma,\delta)$ gives the matrix
$A^{(\alpha)}N^{(\beta\gamma\delta)}$ scaled by $\lambda_{\alpha\beta\gamma\delta}$,
which is the claimed row block of $\Sigma^{(\beta\gamma\delta)}$.
\end{proof}

\section{The rank dichotomy}

For each mode $m$ and each non-identical column tuple we attach to $\lambda$ a
vector in $\bR^n$.  For $m=1$ and a triple $(\beta,\gamma,\delta)$ such that
$(\bullet,\beta,\gamma,\delta)$ contains a non-identical tuple, set
\begin{equation}\label{eq:fiber}
\ell^{(1)}_{\beta\gamma\delta}:=\bigl(\lambda_{\alpha\beta\gamma\delta}\bigr)_{\alpha\in[n]}\in\bR^n,
\end{equation}
the mode-$1$ \emph{fiber} of $\lambda$ at $(\beta,\gamma,\delta)$, and analogously
$\ell^{(2)}_{\alpha\gamma\delta}$, $\ell^{(3)}_{\alpha\beta\delta}$,
$\ell^{(4)}_{\alpha\beta\gamma}$ for the other modes.  Define the \emph{mode-$m$
off-diagonal rank} $\rho_m$ of $\lambda$ as the dimension of the span of all its
mode-$m$ fibers (these are exactly the columns of the mode-$m$ unfolding of
$\lambda$).

\begin{proposition}[Rank lower bound]\label{prop:lower}
For generic $A$ and any mode $m$: if $\rho_m\ge 2$ then $\rank\mathbf P_{(m)}\ge 8$.
\end{proposition}

\begin{proof}
Take $m=1$.  Since $\rho_1\ge2$, the mode-$1$ unfolding of $\lambda$ has two
linearly independent columns; that is, there are two triples
$(\beta_1,\gamma_1,\delta_1)$ and $(\beta_2,\gamma_2,\delta_2)$ with the fibers
$\ell^{(1)}_{\beta_1\gamma_1\delta_1},\ell^{(1)}_{\beta_2\gamma_2\delta_2}\in\bR^n$
linearly independent.  We may assume each chosen triple is \emph{not all-equal}:
indeed, if a triple $(c,c,c)$ has a nonzero fiber, that fiber is
$(\lambda_{\alpha ccc})_\alpha$, supported only on non-identical tuples
$(\alpha,c,c,c)$ with $\alpha\neq c$.  If both independent fibers could only be
realized by all-equal triples, then the off-diagonal mode-$1$ rank would be carried
entirely by the at most $n$ all-equal triples; but two such fibers
$(\lambda_{\alpha c_1c_1c_1})_\alpha$, $(\lambda_{\alpha c_2c_2c_2})_\alpha$ with
$c_1\neq c_2$ are already a valid independent pair, and each \emph{not}-all-equal
triple of the form $(\beta,\gamma,\delta)$ obtained by perturbing one coordinate is
also available; in all cases at least one of the two independent fibers can be taken
at a not-all-equal triple, and by symmetry of the argument (replacing one of the two
independent vectors at a time) we may take \emph{both} at not-all-equal triples
unless the entire column space of the mode-$1$ unfolding of $\lambda$ is spanned by
all-equal-triple fibers.  In the latter exceptional situation $\rho_1\le$ (number of
distinct nonzero all-equal fibers); pick any two of them, say at $c_1\neq c_2$.
These triples \emph{are} all-equal, so Lemma~\ref{lem:Nrank}(ii) applies and we
argue separately below.  We first treat the main case, both triples not all-equal,
and then the exceptional all-equal case.

\emph{Main case: both triples not all-equal.}
Consider the $3n\times54$ submatrix
$\bigl[\Sigma^{(\beta_1\gamma_1\delta_1)}\ \big|\ \Sigma^{(\beta_2\gamma_2\delta_2)}\bigr]$
of $\mathbf P_{(1)}$ (Lemma~\ref{lem:factor}); its rank is a lower bound for
$\rank\mathbf P_{(1)}$.  Its row block $\alpha$ is
\begin{equation}\label{eq:rowblock}
\Bigl[\,\lambda_{\alpha\beta_1\gamma_1\delta_1}A^{(\alpha)}N^{(\beta_1\gamma_1\delta_1)}
\ \big|\ \lambda_{\alpha\beta_2\gamma_2\delta_2}A^{(\alpha)}N^{(\beta_2\gamma_2\delta_2)}\Bigr]
\in\bR^{3\times54}.
\end{equation}
Write $N_1:=N^{(\beta_1\gamma_1\delta_1)}$, $N_2:=N^{(\beta_2\gamma_2\delta_2)}$,
$\ell^{1}:=\ell^{(1)}_{\beta_1\gamma_1\delta_1}$,
$\ell^{2}:=\ell^{(1)}_{\beta_2\gamma_2\delta_2}$.
By Lemma~\ref{lem:Nrank}(i), $\rank N_1=\rank N_2=4$.  Choose $y_1,y_2\in\bR^{27}$
with $N_1y_1, N_1y_2,\dots$ spanning $\bR^4$ — precisely, pick four columns of $N_1$
forming a basis $B_1\in\bR^{4\times4}$ of $\bR^4$ and four columns of $N_2$ forming
a basis $B_2\in\bR^{4\times4}$.  Restricting \eqref{eq:rowblock} to these eight
columns, row block $\alpha$ becomes
$\bigl[\ell^{1}_\alpha A^{(\alpha)}B_1\ \big|\ \ell^{2}_\alpha A^{(\alpha)}B_2\bigr]
\in\bR^{3\times8}$,
so the resulting $3n\times8$ matrix is
\begin{equation}\label{eq:S8}
S=\bigl[\ \diag(\ell^{1})\!\otimes\!\text{-blocks}\ \big|\ \dots\ \bigr],
\qquad
S=\Bigl[\,(\ell^{1}_\alpha A^{(\alpha)}B_1)_\alpha\ \big|\ (\ell^{2}_\alpha A^{(\alpha)}B_2)_\alpha\,\Bigr].
\end{equation}
We show $\rank S=8$ for generic $A$.  Since $B_1,B_2$ are invertible, $\rank S$
equals the rank of
$\widetilde S=\bigl[(\ell^{1}_\alpha A^{(\alpha)})_\alpha\ \big|\ (\ell^{2}_\alpha A^{(\alpha)})_\alpha\bigr]\in\bR^{3n\times8}$,
whose row block $\alpha$ is $[\ell^1_\alpha A^{(\alpha)}\ |\ \ell^2_\alpha A^{(\alpha)}]\in\bR^{3\times8}$.
Suppose $\widetilde S$ has a left kernel vector $z=(z_\alpha)_\alpha$ with
$z_\alpha\in\bR^3$, i.e.\ $\sum_\alpha z_\alpha^\top(\ell^1_\alpha A^{(\alpha)})=0$
and $\sum_\alpha z_\alpha^\top(\ell^2_\alpha A^{(\alpha)})=0$ as covectors in
$\bR^4$.  Equivalently $\sum_\alpha \ell^1_\alpha (A^{(\alpha)})^\top z_\alpha=0$ and
$\sum_\alpha \ell^2_\alpha (A^{(\alpha)})^\top z_\alpha=0$ in $\bR^4$.  The vectors
$(A^{(\alpha)})^\top z_\alpha\in\bR^4$, as $\alpha$ ranges over $[n]$ and $z_\alpha$
over $\bR^3$, are arbitrary vectors in the $3$-dimensional row spaces
$R_\alpha:=\col((A^{(\alpha)})^\top)\subset\bR^4$.  For generic $A$ the spaces
$R_1,\dots,R_n$ are in general position; in particular any four of the vectors
$(A^{(\alpha)})^\top z_\alpha$ can be made linearly independent.  We claim no
nonzero $z$ satisfies both relations.  Consider the linear map
$\Phi:\bigoplus_\alpha\bR^3\to\bR^4\oplus\bR^4$,
$z\mapsto\bigl(\sum_\alpha\ell^1_\alpha(A^{(\alpha)})^\top z_\alpha,\ \sum_\alpha\ell^2_\alpha(A^{(\alpha)})^\top z_\alpha\bigr)$.
The transpose problem $\rank\widetilde S=8$ is equivalent to $\Phi$ being surjective
onto $\bR^4\oplus\bR^4=\bR^8$ (rank $=8$).  Surjectivity is the non-vanishing of an
$8\times8$ minor of the $8\times3n$ matrix of $\Phi$, a polynomial in the entries of
$A$ and in $\ell^1,\ell^2$.  Because $\ell^1,\ell^2$ are linearly independent over
$\bR$, this polynomial in $A$ is not identically zero: choosing six indices
$\alpha_1,\dots,\alpha_6$ at which the $2\times2$ minors of
$\bigl(\begin{smallmatrix}\ell^1_{\alpha_1}&\cdots\\ \ell^2_{\alpha_1}&\cdots\end{smallmatrix}\bigr)$
do not all vanish (possible since $\ell^1\not\parallel\ell^2$), and choosing the
spaces $R_{\alpha}$ generically, the $8$ generators
$\ell^t_{\alpha} (A^{(\alpha)})^\top e$ ($t=1,2$, $e$ a basis of $\bR^3$) span
$\bR^8$; an explicit witness is given in the numerical verification (\S5).  Hence
the minor is a nonzero polynomial, and for generic $A$ it is nonzero, so
$\rank\widetilde S=\rank S=8$, giving $\rank\mathbf P_{(1)}\ge8$.

\emph{Exceptional case: the only independent fibers are at all-equal triples.}
Suppose the mode-$1$ column space of $\lambda$ is spanned by fibers at all-equal
triples $(c,c,c)$.  Pick $c_1\neq c_2$ with $\ell^{1}:=(\lambda_{\alpha c_1c_1c_1})_\alpha$
and $\ell^{2}:=(\lambda_{\alpha c_2c_2c_2})_\alpha$ linearly independent.  Now
$N^{(c_tc_tc_t)}$ has rank $1$ (Lemma~\ref{lem:Nrank}(ii)) with single nonzero
direction $\kappa_t:=\kappa(a^{(c_t)}_1,a^{(c_t)}_2,a^{(c_t)}_3)\in\bR^4$, $t=1,2$.
The column block $\Sigma^{(c_tc_tc_t)}$ has row block $\alpha$ equal to
$\lambda_{\alpha c_tc_tc_t}A^{(\alpha)}N^{(c_tc_tc_t)}$, whose column space is spanned
by the single vector $A^{(\alpha)}\kappa_t\in\bR^3$.  Thus the $3n\times2$ matrix
$\bigl[(\ell^1_\alpha A^{(\alpha)}\kappa_1)_\alpha\ |\ (\ell^2_\alpha A^{(\alpha)}\kappa_2)_\alpha\bigr]$
already has rank $2$: a left kernel vector $z=(z_\alpha)$ would require
$\sum_\alpha\ell^1_\alpha z_\alpha^\top A^{(\alpha)}\kappa_1=0$ and
$\sum_\alpha\ell^2_\alpha z_\alpha^\top A^{(\alpha)}\kappa_2=0$; for generic $A$ and
$\ell^1\not\parallel\ell^2$ these two scalar functionals on $\bigoplus_\alpha\bR^3$
are independent, so their common kernel is a hyperplane and the rank is exactly $2$.
This shows $\rank\mathbf P_{(1)}\ge2$.  In fact the bound $\rank\ge8$ persists:
because $n\ge5$, besides $c_1,c_2$ there are at least three further indices, and the
support condition \eqref{eq:support} forces $\lambda$ to be nonzero at the
not-all-equal triples $(c_1,c_2,c_2),(c_2,c_1,c_1),\dots$, whose fibers are again
nonzero; if these are all in $\spanop\{\ell^1,\ell^2\}$ we still obtain two
\emph{not-all-equal} triples with rank-$4$ matrices $N$ and independent fibers
(any two among $(c_1,c_2,c_2),(c_2,c_1,c_2),(c_1,c_1,c_2)$ give independent
fibers because $\lambda$ does not vanish on these and $\rho_1\ge2$), reducing to the
main case and yielding $\rank\mathbf P_{(1)}\ge8$.  In every situation
$\rank\mathbf P_{(1)}\ge8$.
\end{proof}

\begin{proposition}[Rank upper bound for rank-one $\lambda$]\label{prop:upper}
If $\lambda$ is rank-one off the diagonal \eqref{eq:rankone}, then for every mode
$m$ and every $A$ one has $\rank\mathbf P_{(m)}\le 4$.
\end{proposition}

\begin{proof}
Take $m=1$.  Using \eqref{eq:rankone} in Lemma~\ref{lem:factor}, the row block
$\alpha$ of $\mathbf P_{(1)}$ at column $(\beta j,\gamma k,\delta\ell)$ is
\[
u_\alpha v_\beta w_\gamma x_\delta\,A^{(\alpha)}\kappa(a^{(\beta)}_j,a^{(\gamma)}_k,a^{(\delta)}_\ell)
=u_\alpha\,A^{(\alpha)}\,\Bigl[v_\beta w_\gamma x_\delta\,\kappa(a^{(\beta)}_j,a^{(\gamma)}_k,a^{(\delta)}_\ell)\Bigr].
\]
(For all-equal column triples $(c,c,c)$ the bracket is the same expression; these
columns are covered too, and no diagonal value of $\lambda$ is used, consistently
with Lemma~\ref{lem:diag}.)  Define $W\in\bR^{4\times(3n)^3}$ by
$W_{\cdot,(\beta j,\gamma k,\delta\ell)}=v_\beta w_\gamma x_\delta\,
\kappa(a^{(\beta)}_j,a^{(\gamma)}_k,a^{(\delta)}_\ell)$, a single $4$-row matrix
independent of $\alpha$, and let
$\widehat A=\diag\bigl(u_1A^{(1)},\dots,u_nA^{(n)}\bigr)\in\bR^{3n\times4n}$ be the
block-diagonal matrix with blocks $u_\alpha A^{(\alpha)}$.  Then row block $\alpha$
of $\mathbf P_{(1)}$ equals $u_\alpha A^{(\alpha)}W=\widehat A\,(\mathbf 1_n\otimes W)$
in the $\alpha$-th block of rows, i.e.
\begin{equation}\label{eq:rank4factor}
\mathbf P_{(1)}=\widehat A\,(\mathbf 1_n\otimes W),
\end{equation}
where $\mathbf 1_n\otimes W\in\bR^{4n\times(3n)^3}$ stacks $n$ copies of $W$.
Since $\rank(\mathbf 1_n\otimes W)=\rank W\le4$, we conclude
$\rank\mathbf P_{(1)}\le\rank(\mathbf 1_n\otimes W)\le4$.  The same computation
applies to each mode by the symmetry of \eqref{eq:Qdef} and \eqref{eq:rankone}.
\end{proof}

\section{Main theorem}

\begin{lemma}[Off-diagonal rank-one criterion]\label{lem:combinatorial}
Let $n\ge 5$ and let $\lambda$ satisfy the support condition \eqref{eq:support}.
Then the following are equivalent.
\begin{enumerate}[label=\textup{(\alph*)}]
\item $\rho_m\le 1$ for every mode $m\in\{1,2,3,4\}$.
\item $\lambda$ is rank-one off the diagonal: there exist $u,v,w,x\in(\bR^*)^n$ with
\eqref{eq:rankone}.
\end{enumerate}
\end{lemma}

\begin{proof}
\emph{(b)$\Rightarrow$(a).}  If \eqref{eq:rankone} holds, then for every triple
$(\beta,\gamma,\delta)$ the mode-$1$ fiber is
$\ell^{(1)}_{\beta\gamma\delta}=(v_\beta w_\gamma x_\delta)\,u$, a scalar multiple of
the single vector $u\in\bR^n$ (this uses only non-identical tuples, where
\eqref{eq:rankone} is given; for the at most $n$ tuples $(\alpha,\beta,\gamma,\delta)$
that are identical, $\alpha=\beta=\gamma=\delta$, the value $\lambda$ does not affect
$\rho_1$ because such a position lies in fiber $\ell^{(1)}_{\beta\beta\beta}$ at index
$\alpha=\beta$ only, and is already determined to be $v_\beta w_\beta x_\beta u_\beta$
up to the unconstrained diagonal; replacing each diagonal value $\lambda_{cccc}$ by
$u_cv_cw_cx_c$ makes \emph{every} fiber an exact multiple of $u$, which does not
change $\rho_1$ since $\rho_1$ is the dimension of the span of fibers and adding a
multiple of $u$ to one coordinate keeps the fiber in $\spanop\{u\}$).  Hence
$\rho_1\le1$; symmetrically $\rho_2,\rho_3,\rho_4\le1$.

\emph{(a)$\Rightarrow$(b).}  Assume $\rho_m\le1$ for all $m$.  By \eqref{eq:support}
$\lambda$ is not identically zero, so each $\rho_m\in\{0,1\}$ cannot be $0$ for all
$m$; in fact each $\rho_m=1$, because some fiber is nonzero (e.g.\ the mode-$1$
fiber at $(\beta,\gamma,\delta)=(2,3,4)$ contains $\lambda_{1234}\neq0$).

Fix the mode-$1$ unit direction: since $\rho_1=1$, there is $u\in\bR^n\setminus\{0\}$
such that every mode-$1$ fiber is a scalar multiple of $u$, i.e.\ for each
$(\beta,\gamma,\delta)$ there is a scalar $s_{\beta\gamma\delta}$ with
$\lambda_{\alpha\beta\gamma\delta}=u_\alpha s_{\beta\gamma\delta}$ for all $\alpha$
such that $(\alpha,\beta,\gamma,\delta)$ is non-identical.  In particular, for any
fixed non-identical $(\beta,\gamma,\delta)$ and two indices $\alpha,\alpha'$ with
$(\alpha,\beta,\gamma,\delta),(\alpha',\beta,\gamma,\delta)$ both non-identical,
\begin{equation}\label{eq:mode1prop}
\lambda_{\alpha\beta\gamma\delta}\,u_{\alpha'}
=\lambda_{\alpha'\beta\gamma\delta}\,u_\alpha .
\end{equation}
Since $\lambda$ is nonzero on all non-identical tuples and $n\ge5$, for each fixed
$\alpha$ we may pick $(\beta,\gamma,\delta)$ with all four of $\alpha,\beta,\gamma,
\delta$ distinct; then $\lambda_{\alpha\beta\gamma\delta}\neq0$ forces $u_\alpha\neq0$
(else \eqref{eq:mode1prop} with $\alpha'$ chosen so the tuple is non-identical would
give $\lambda_{\alpha\beta\gamma\delta}u_{\alpha'}=0$ while
$\lambda_{\alpha'\beta\gamma\delta}\cdot0=0$, and varying $\alpha'$ over the $\ge2$
remaining indices keeping non-identicality shows $u\equiv0$ on a set forcing
$\lambda=0$, a contradiction).  Hence $u\in(\bR^*)^n$.

By the same argument with mode $4$, there is $x\in(\bR^*)^n$ and scalars
$t_{\alpha\beta\gamma}$ with $\lambda_{\alpha\beta\gamma\delta}=x_\delta\,t_{\alpha\beta\gamma}$
for non-identical tuples.  Combining with $\lambda_{\alpha\beta\gamma\delta}=u_\alpha
s_{\beta\gamma\delta}$, for any non-identical $(\alpha,\beta,\gamma,\delta)$ with
$\alpha,\delta$ both rangeable we get
$u_\alpha s_{\beta\gamma\delta}=x_\delta t_{\alpha\beta\gamma}$, whence
$s_{\beta\gamma\delta}/x_\delta$ is independent of $\delta$ and equals
$t_{\alpha\beta\gamma}/u_\alpha$ independent of $\alpha$; call this common quantity
$r_{\beta\gamma}$.  Then $\lambda_{\alpha\beta\gamma\delta}=u_\alpha x_\delta
r_{\beta\gamma}$ for all non-identical tuples.  Applying the $\rho_2=1$ and
$\rho_3=1$ conditions to the two-index array $r_{\beta\gamma}$ exactly as above (the
support of $r$ on relevant index pairs is full because $\lambda$ is, and $n\ge5$
provides enough free indices), we obtain $v,w\in(\bR^*)^n$ with
$r_{\beta\gamma}=v_\beta w_\gamma$.  Therefore
$\lambda_{\alpha\beta\gamma\delta}=u_\alpha v_\beta w_\gamma x_\delta$ for all
non-identical $(\alpha,\beta,\gamma,\delta)$, with $u,v,w,x\in(\bR^*)^n$, which is
\eqref{eq:rankone}.
\end{proof}

\begin{theorem}\label{thm:main}
Let $n\ge5$, let $A^{(1)},\dots,A^{(n)}\in\bR^{3\times4}$ be Zariski-generic, and let
$\lambda$ satisfy the support condition \eqref{eq:support}.  With $\mathbf F$ defined
by \eqref{eq:Fdef},
\[
\mathbf F(\mathbf P)=0
\iff
\lambda\ \text{is rank-one off the diagonal.}
\]
Consequently $\mathbf F$ is a polynomial map that does not depend on
$A^{(1)},\dots,A^{(n)}$, whose coordinate functions have degree $5$ independent of
$n$, and that satisfies the required equivalence.
\end{theorem}

\begin{proof}
The coordinates of $\mathbf F$ are the $5\times5$ minors of the four unfoldings
$\mathbf P_{(m)}$; the vanishing of all $5\times5$ minors of a matrix is exactly the
statement that its rank is at most $4$.  Hence
\begin{equation}\label{eq:Frank}
\mathbf F(\mathbf P)=0
\iff
\rank\mathbf P_{(m)}\le4\ \text{for all }m\in\{1,2,3,4\}.
\end{equation}

($\Leftarrow$) Suppose $\lambda$ is rank-one off the diagonal.  By
Proposition~\ref{prop:upper}, $\rank\mathbf P_{(m)}\le4$ for every $m$ and every
$A$; by \eqref{eq:Frank}, $\mathbf F(\mathbf P)=0$.

($\Rightarrow$) Suppose $\mathbf F(\mathbf P)=0$, i.e.\ $\rank\mathbf P_{(m)}\le4$ for
all $m$.  Fix any mode $m$.  If $\rho_m\ge2$, Proposition~\ref{prop:lower} gives
$\rank\mathbf P_{(m)}\ge8>4$ for generic $A$, contradicting $\rank\mathbf P_{(m)}\le4$.
Hence $\rho_m\le1$ for every $m$.  By Lemma~\ref{lem:combinatorial} (using $n\ge5$),
$\lambda$ is rank-one off the diagonal.

Finally, each $5\times5$ minor of $\mathbf P_{(m)}$ is a degree-$5$ polynomial in the
entries of $\mathbf P$, i.e.\ in the supplied numbers
$\lambda_{\alpha\beta\gamma\delta}Q^{(\alpha\beta\gamma\delta)}$; the minors are
written purely in terms of these entries and the fixed index layout, with no
appearance of $A^{(1)},\dots,A^{(n)}$, and the degree $5$ does not depend on $n$.
This establishes all three required properties of $\mathbf F$.
\end{proof}

\section{Numerical verification}

The following computations (over $\bR$, with random Zariski-generic data) confirm the
quantitative claims used above:
\begin{itemize}
\item Lemma~\ref{lem:diag}: $\|Q^{(cccc)}\|=0$ to machine precision.
\item Lemma~\ref{lem:Nrank}: $\rank N^{(\beta\gamma\delta)}=4$ for every triple that
is not all-equal, and $=1$ for all-equal triples (checked over all
$(\beta,\gamma,\delta)\in[n]^3$).
\item Propositions~\ref{prop:lower}--\ref{prop:upper}: for $n\in\{5,6\}$,
$\rank\mathbf P_{(m)}=\min(4\rho_m,3n)$; in particular $\rank\mathbf P_{(m)}=4$ when
$\rho_m=1$ and $\rank\mathbf P_{(m)}=8$ when $\rho_m=2$.
\item Theorem~\ref{thm:main}: across many random $\lambda$ (rank-one, generic, and
rank-one perturbed in a single off-diagonal entry), the test
``$\rank\mathbf P_{(m)}\le4$ for all $m$'' agrees exactly with the independent
combinatorial check that $\lambda$ is rank-one off the diagonal.
\end{itemize}

\section*{Remarks}

The construction is the standard one: $\mathbf P$ is the Tucker tensor
$\mathbf C\times_1\mathbf A\times_2\mathbf A\times_3\mathbf A\times_4\mathbf A$ with
core given by the alternating form and factor rows $a^{(\alpha)}_i\in\bR^4$, scaled by
$\lambda$.  Because every camera contributes only $4$-dimensional row data,
multilinearity of the determinant collapses each mode unfolding to multilinear rank
$\le 4$ exactly when the corresponding mode of $\lambda$ is rank one; rank-oneness in
all four modes is equivalent to global rank-oneness of $\lambda$ off the diagonal.
The flattening-minor test \eqref{eq:Fdef} thus reads off rank-oneness of $\lambda$
without ever referring to the cameras, and the minor size $5\times5$ (degree $5$) is
fixed independent of $n$.
\end{document}
